\section{Conclusion}

We carried out a quantitative XAI assessment of a ResNet-50 model on the public UIdataGB dataset for nine-class gallbladder ultrasound classification, and along the way found that our own first-pass leakage correction was itself incomplete. An initial split, de-duplicated only by single-orientation perceptual hashing, gave a 93.05\% headline accuracy; a closer audit found that near-identical frames from the same imaging session, and rotation-augmented duplicates the dataset's release pipeline introduces, still crossed the train/validation boundary under that split. A corrected, cluster-aware 70/20/10 split---with a genuine test set touched exactly once---puts real, honestly-measured accuracy at 72.56\%~$\pm$~2.35\% across four independently seeded models. We treat that lower, corrected number as the paper's real result, and the 93.05\% figure as a demonstration of how much an incomplete leakage correction can inflate apparent performance.

The main XAI comparison---Grad-CAM, Grad-CAM++, gradient Saliency Maps, Eigen-CAM, and Score-CAM, evaluated via a cosine-based consistency metric $O^c$ and a reliability suite covering faithfulness, stability, a model-randomisation sanity check, calibration, and a shortcut audit---was run on the higher-accuracy, less-corrected split. Grad-CAM++ and Score-CAM produced activation maps that were sharper, more class-discriminative, and substantially more consistent across replicates than the other methods ($\bar{O}^c=0.944$ and $0.937$ versus 0.769, 0.242, and 0.681), though that consistency edge did not carry over to faithfulness, where all three CAM-family methods were similar and plain Grad-CAM had the largest deletion/insertion gap: the choice of method depends on the metric you care about. Those original consistency replicates, however, were trained on disjoint data subsets rather than identical data with only the seed varied, which conflates data-composition effects with true reproducibility. We reran the consistency check under a corrected same-data design on the leakage-fixed split, using Grad-CAM as a representative method, alongside a null control that the original metric had never been tested against: cosine similarity between untrained models, and between synthetic random heatmaps of the same shape. The corrected consistency score clears both floors by roughly 2.5$\times$, so the finding survives this test. We then extended the corrected design to the complete battery---consistency (all five methods), faithfulness (the three feature-map-space CAM methods), stability, the sanity check, and the shortcut audit (Grad-CAM++ as the representative method), all re-run on the Stage~2 checkpoints---and found the paper's central method-comparison conclusions hold in direction: Grad-CAM++ and Score-CAM remain the most consistent methods, and no method dominates on every axis, exactly as under Stage~1. What changes with the corrected split is calibration quality, which degrades in step with the lower, honest accuracy figure, not independently of it.

The broader reliability battery showed that the explanations were faithful to the model's decisions, held up when brightness was altered, relied on the learned weights, and were not driven by margin-based shortcuts specifically. That last finding, however, is narrower than it first appears: the margin audit was never able to detect a batch-correlated shortcut by design, and a targeted follow-up audit confirms one exists. A batch-ID diagnostic, re-verified so no duplicate cluster crosses its own train/validation split, still reaches 82.88\% accuracy (8.3$\times$ chance) at telling acquisition batches apart within a single disease class, and a linear probe on the disease classifier's own features recovers batch identity at 74.75\% (4.5$\times$ chance)---direct evidence the shortcut lives inside the very classifier this paper's explanations describe, not just in a separately-trained proxy. A noise-robustness sweep adds a second, more severe caution: classification accuracy on the corrected model falls from 90.0\% clean to 20.6\%, close to chance, under a still-subtle noise perturbation, a far larger effect than anything in the XAI comparison itself. Failure-case analysis further revealed that a plausible heatmap alone is not enough to flag high-confidence misclassifications. This combination of a corrected, honestly-reported accuracy figure, a stress-tested XAI evaluation, and these follow-up audits is the paper's core contribution.

These findings make sense only if the data are truly held-out. The corrected split ensures this at both the perceptual-hash and imaging-session level, with duplicate clusters never crossing a train/validation/test boundary (see Sections~\ref{sec:methods}.1 and \ref{sec:stage2}).

The study has several limitations, all of which are laid out in Section~\ref{sec:limitations}. First, there is no clinician-annotated localization to back the visual explanations. Second, the full five-method XAI battery and the failure-case visual analysis have both now been re-run on the corrected split (Sections~\ref{sec:stage2-failure-cases} and~\ref{sec:stage2-full-battery}), closing what was previously an open Stage~1/Stage~2 method gap. Third, faithfulness and stability were tested only with a bounded set of perturbations. Fourth, the shortcut audit's margin-crop leg tests only margin-tied shortcuts; the batch-correlated shortcut it could not detect is now confirmed across all nine disease classes (Section~\ref{sec:followup-audits}), though the full nine-class probe uses an image-level rather than cluster-aware per-batch split, so its 100\% per-class figures likely include some of the same near-duplicate-frame inflation the single-class batch-ID diagnostic's cluster-aware re-verification partly corrected for. Fifth, the dataset comes from multiple hospitals and scanners but lacks labels for that provenance (see Section~\ref{sec:future-research}). Sixth, the pHash threshold sensitivity now includes real per-threshold accuracy (Section~\ref{sec:followup-audits}), but the noise-robustness sweep remains partial: it does not retrain a classifier across the full noise range tested, so its conclusions describe severity, not an exhaustive characterisation.
